\subsection{1. Đặt vấn đề}

\indent\indent Trong kỷ nguyên Web3, các giá trị của mạng xã hội phi tập trung (SocialFi), tiêu biểu là Lens Protocol, ngày càng được khẳng định nhờ cơ chế trao quyền sở hữu dữ liệu và khả năng tương tác sinh lời cho người dùng. việc xây dựng một hệ sinh thái an toàn và minh bạch không chỉ giúp bảo vệ lợi ích cộng đồng mà còn mở ra cơ hội phát triển cho nền kinh tế số, quản trị phi tập trung và các chương trình tặng thưởng mã thông báo. Tuy nhiên, tính mở và các ưu đãi kinh tế đặc thù của Web3 cũng vô tình trở thành mục tiêu cho các cuộc tấn công giả mạo (Sybil attacks), nơi những cá nhân hoặc tổ chức tạo lập hàng loạt tài khoản ảo nhằm thao túng thuật toán, chiếm quyền kiểm soát quản trị và trục lợi bất chính.\vspace{0.3cm}

Đặc điểm của dữ liệu trên mạng xã hội phi tập trung là tính đa quan hệ và sự phụ thuộc ràng buộc cao. Mỗi thực thể người dùng thường tồn tại đồng thời dưới nhiều tầng dữ liệu: thông tin hồ sơ NFT, hành vi xã hội (theo dõi, phản hồi), tương tác nội dung (thu thập, chia sẻ), và đặc biệt là các mối liên kết tài chính trên chuỗi khối như phí giao dịch hay quyền đồng sở hữu ví. Bên cạnh đó, các mạng lưới tài khoản giả mạo ngày càng tinh vi trong việc ngụy trang metadata để mô phỏng người dùng thật, làm tăng độ khó trong việc phân lớp nếu chỉ dựa vào một phương thức dữ liệu đơn lẻ hoặc các cấu trúc đồ thị truyền thống. Những hạn chế này khiến các phương pháp học máy cơ bản như MLP hay CNN đơn phương thức bỏ sót các mối liên hệ ẩn và các ràng buộc cấu trúc quan trọng trong không gian đồ thị siêu đặc.\vspace{0.3cm}

Trước những thách thức trên, đề tài đề xuất hệ thống SybilSignal, ứng dụng mô hình trí tuệ đồ thị nhận thức ràng buộc (Constraint-Aware) kết hợp mạng nơ-ron đồ thị (GNN) và học chuyển giao. Với cấu trúc đồ thị đa tầng, hệ thống sẽ khai thác đồng thời đặc trưng của từng tài khoản và các ràng buộc tài chính không thể chối bỏ trên chuỗi khối, cho phép mô hình hiểu sâu các ngữ cảnh hành vi tự nhiên cũng như các cấu trúc bầy đàn của trang trại bot. Tiếp đó, kiến trúc phân lớp tách rời (Decoupled ML) được sử dụng để tối ưu hóa không gian nhúng từ bộ mã hóa tự động đồ thị (GAE) và đưa vào các bộ phân loại có tính giải thích cao như Random Forest. Sự kết hợp này kỳ vọng nâng cao hiệu quả phát hiện trong điều kiện thiếu nhãn dữ liệu, vượt trội hơn các mô hình truyền thống, đồng thời cung cấp khả năng diễn giải minh bạch, góp phần bảo vệ tính an toàn và bền vững cho các nền tảng mạng xã hội phi tập trung thế hệ mới.